\section{模型的评价与推广}

\subsection{模型的评价}

\textbf{模型的优点：}
\begin{enumerate}
\item 四问采用统一的热湿迁移框架，依次引入状态相关物性、全域首达判据与材料收缩，能够衔接短期预热和长期干燥预测。
\item 有限体积法直接平衡控制体通量，全隐式与 Picard 迭代适应非线性物性；网格、步长、二维及交叉求解检验支持数值结果的可靠性。
\item 全域阈值判据识别中心控制位置，同物性固定半径对照分离几何收缩效应，便于解释干燥时间差异。
\end{enumerate}

\textbf{模型的局限：}
\begin{enumerate}
\item 中截面近似获得同物性、同边界二维对照的支持，但材料各向异性及蒸发潜热的影响尚未量化；既有数值验证不能替代对这些物理简化的检验。
\item 环境水分采用等效映射，交换系数沿用题给值，长期边界取常值延拓；实际应用需标定浓度转换与表面交换关系，并将后续工况变化纳入预测范围。
\item 收缩模型以给定半径历史和均匀径向缩放闭合内部运动。水分守恒和坐标实现对照支持这一闭合的数值自洽性，实际局部变形、长度变化及完整能量平衡仍需另行验证。
\end{enumerate}

\subsection{模型的推广}

结合少量中心温湿测量，可标定扩散与表面交换参数，利用模型估计内部达标时刻，并进一步加入能耗、品质和设备约束以优化干燥温度策略。

对短柱或端面交换显著的材料，可采用二维轴对称模型；获得潜热和吸湿数据后，可扩展为 Luikov 型热湿耦合模型\cite{luikov1975}。经物性与边界重新标定，该框架可用于果蔬、木材等多孔材料的干燥预测。
